\chapter{Развилка пятнадцати и шестнадцати состояний и степень нейтринного сектора}

Обратный аудит связности обнаружил, что носитель Мориты
$M_{20\times15}(\mathbb C)$ был построен без правого нейтрино, тогда как
полный перенормируемый дираковский граф с правым нейтрино содержит четыре
рёбра Юкавы. До построения физического бимодуля одноформ необходимо решить,
является ли это противоречием или разделением операторов разных степеней.

\section{Антикруговая сверка}

Эта глава не вводит новый потенциал и не подбирает нейтринную массу. Она
сопоставляет четыре уже полученных результата:
\begin{enumerate}
 \item точное наблюдаемое чтение
 $H_{15}=Q_L\oplus L_L\oplus u_R\oplus d_R\oplus e_R$;
 \item связывающий родитель
 $M_{20\times15}(\mathbb C)\subset M_{35}(\mathbb C)$;
 \item четыре перенормируемых дираковских ребра при наличии $\nu_R$;
 \item условный майорановский дефект, топология которого согласована, но
 динамическое возникновение ещё не выведено.
\end{enumerate}
В литературе отсутствие $\nu_R$ означает, что первая калибровочно
инвариантная масса нейтрино возникает как оператор Вайнберга размерности
пять. Поэтому различие между $H_{15}$ и $H_{16}$ одновременно является
различием полевого состава и степени оператора.

\section{Две несовпадающие архитектуры}

Для текущего наблюдаемого пакета
\begin{equation}
 15=6+2+3+3+1
\end{equation}
имеем
\begin{equation}
 E_{15}=M_{20\times15}(\mathbb C),\qquad
 \mathcal L(E_{15})=M_{35}(\mathbb C),
\end{equation}
а единственный нормированный след даёт угловые веса
\begin{equation}
 \tau_{35}(p_{20})=\frac47,
 \qquad
 \tau_{35}(p_{15})=\frac37.
\end{equation}

Добавление одного правого нейтрино создаёт не тот же объект с ещё одним
ребром, а другую архитектуру:
\begin{equation}
 E_{16}=M_{20\times16}(\mathbb C),\qquad
 \mathcal L(E_{16})=M_{36}(\mathbb C),
\end{equation}
\begin{equation}
 \tau_{36}(p_{20})=\frac59,
 \qquad
 \tau_{36}(p_{16})=\frac49.
\end{equation}
Размерность переходного носителя меняется с $300$ на $320$, а ранее
выведенные веса $4/7$ и $3/7$ исчезают. Следовательно, $H_{16}$ требует
повторной проверки общего следа, коммутирующего квадрата, углового чтения и
всех последующих нормировок. Его нельзя незаметно вставить в родитель
$M_{35}$.

\section{Разделение первой и высшей степеней}

На $H_{15}$ один хиггсовский дублет допускает ровно три
перенормируемых заряженных ребра
\begin{equation}
 \overline Q_L\widetilde H u_R,
 \qquad
 \overline Q_L H d_R,
 \qquad
 \overline L_L H e_R.
\end{equation}
Ребро $\overline L_L\widetilde H\nu_R$ отсутствует не вследствие
проекции, а вследствие отсутствия самого правого модуля. Поэтому
физический модуль одноформ на текущем носителе обязан сначала отвечать
только за $u,d,e$.

Нейтринный сектор при этом не объявляется нулевым. Он переносится в
отдельную задачу высшей степени:
\begin{equation}
 \frac{c_{ij}}{\Lambda}
 (\overline{L_i^c}\,\widetilde H^*)
 (\widetilde H^\dagger L_j),
\end{equation}
либо в уже исследованную майорановскую дефектную ветвь. Но проект ранее
доказал только условный топологический результат
\begin{equation}
 \pi\text{-поток}+\Phi\ne0
 \Longrightarrow \operatorname{wind}\Phi=1,
\end{equation}
а не происхождение конденсата $\Phi$. Поэтому разделение степеней не
закрывает нейтринную динамику и не позволяет подставлять её коэффициенты.

\begin{theorem}[Минимальная заморозка наблюдаемого угла]
Если сохраняются уже выведенные носитель $M_{20\times15}(\mathbb C)$,
связывающая алгебра $M_{35}(\mathbb C)$ и угловые веса $4/7,3/7$, то
ближайший физический модуль одноформ содержит ровно три заряженных типа
Юкавы $u,d,e$. Нейтринная масса принадлежит отдельному оператору высшей
степени. Переход к четырём дираковским рёбрам требует новой архитектуры
$M_{36}(\mathbb C)$ и повторения нормировочных гейтов.
\end{theorem}

\begin{nogocriterion}[Запрет скрытого шестнадцатого состояния]
Нельзя использовать $Y_\nu$ или дираковское ребро $L_L\to\nu_R$ внутри
носителя $H_{15}$ и одновременно сохранять размерность $300$, алгебру
$M_{35}$ и веса $4/7,3/7$. Такой расчёт смешивал бы две разные
архитектуры.
\end{nogocriterion}

\begin{protocol}
\begin{enumerate}
 \item точное совпадение $H_{15}$ с текущим родителем $M_{35}$:
 \textbf{получено};
 \item число заряженных перенормируемых рёбер на $H_{15}$:
 \textbf{три};
 \item дираковское нейтринное ребро на $H_{15}$:
 \textbf{отсутствует};
 \item безвредность перехода к $H_{16}$:
 \textbf{опровергнута изменением носителя и весов следа};
 \item нейтринный оператор высшей степени:
 \textbf{структурно допустим, динамически не выведен};
 \item архитектура следующего гейта:
 \textbf{$H_{15}$ и три заряженных одноформенных ребра};
 \item физическое замыкание:
 \textbf{не достигнуто}.
\end{enumerate}
\end{protocol}

Следующий гейт должен построить физический бимодуль одноформ на $H_{15}$,
проверить его проективность, дифференцируемость, KO6 и грассманову
связность ранг-один проектора. До этого новый потенциал не вводится.

Машинный сертификат сохранён в
\path{s2t_v5_h15_neutrino_degree_split_gate.py} и
\path{s2t_v5_h15_neutrino_degree_split_gate_results.json}.